\chapter*{前言}
\addcontentsline{toc}{chapter}{前言}

当你第一次在命令行敲下 \texttt{torchrun --nproc\_per\_node=8 train.py}，看到八张 GPU 同时开始吐出 loss 曲线时，也许会感到一种强烈的兴奋：大模型训练，终于跑起来了。但兴奋过后，问题很快浮现——显存不够、梯度 NaN、通信效率低下、推理延迟抖动、KV Cache 占满、集群调度混乱……你发现，真正让大模型跑起来并且跑得好的，不只是写好一行 \texttt{model.forward()}，而是一整套从数学、硬件、通信到调度都紧密耦合的系统工程。

这正是 AI Infrastructure 所处的位置。它不是某个单一框架的配置手册，也不是 GPU 编程的入门教程，更不是模型算法的论文综述。AI Infra 是那条把数学公式落到张量算子、把算子落到 CUDA kernel、把 kernel 落到 GPU 硬件、把多 GPU 通信落到网络拓扑、把整个训练和推理流程落到集群调度的长链条。这条链条上的每一环，都可能成为瓶颈；而理解整条链条，才是从"能用"走向"好用"的关键。

\section*{为什么写这本书}

近年来，大语言模型的发展速度惊人。GPT、LLaMA、Qwen、DeepSeek 等模型不断刷新人们对 AI 能力的认知，开源生态也让越来越多的工程师有机会接触、微调、部署甚至训练这些模型。然而，与模型能力的快速发展相比，AI Infrastructure 的知识体系仍然分散在论文、博客、代码仓库和内部文档中，缺乏一份系统性的、从基础到前沿的完整指南。

很多工程师的典型学习路径是这样的：先从 PyTorch 入门，学会写训练脚本；然后遇到显存不够，开始查 ZeRO 和 FSDP 的配置；再遇到训练不稳定，开始搜索 loss spike 和 gradient overflow 的帖子；接着尝试部署推理服务，发现 KV Cache 和 continuous batching 的概念从未系统学过；最终面对集群调度和网络拓扑，才意识到自己对 RDMA 和 NVLink 几乎一无所知。

这条路径的特点是\textbf{碎片化}——每个问题都是单独搜索、单独解决，缺乏全局视角。工程师可能在某个具体问题上积累了经验，却难以把不同层级的知识串联起来。一个熟悉 NCCL 调试的人，可能不理解 FlashAttention 为什么是 IO 优化而非近似算法；一个熟悉 vLLM 配置的人，可能不知道 PagedAttention 的 block table 与操作系统页表之间的类比；一个熟悉 tensor parallel 的人，可能不清楚 pipeline parallel 的 bubble 与 micro-batch 数量的精确关系。

本书的目标，正是填补这条知识链条上的空白，帮助读者建立\textbf{从数学公式到硬件成本的完整心智模型}。

\section*{这本书的核心理念}

贯穿全书的一条主线是：

\begin{center}
\textbf{数学对象 $\rightarrow$ 张量算子 $\rightarrow$ CUDA kernel $\rightarrow$ GPU 硬件 $\rightarrow$ 多 GPU 通信 $\rightarrow$ 集群调度}
\end{center}

我们不孤立地讲任何一个环节。每一章都会追问：这个概念在上面是什么数学形式，在下面是什么硬件成本，在旁边是什么系统 trade-off。

具体来说，本书坚持三个原则：

第一，\textbf{不只讲"怎么做"，更要讲"为什么"。}很多框架文档会告诉你 ZeRO Stage 3 怎么配置，但不会告诉你为什么 optimizer state 的显存是 $16N$ bytes，为什么 reduce-scatter 比 all-reduce 更适合梯度分片，为什么 ZeRO Stage 3 在某些场景下反而比 Stage 2 更慢。本书在每一个关键决策点，都会从数学和系统两个层面解释因果。

第二，\textbf{不回避底层，也不脱离上层。}我们会深入 CUDA 的 grid/block/thread、shared memory 和 Tensor Core，因为不理解这些就无法解释 FlashAttention 的 tiling 策略和 PagedAttention 的访存模式。但我们不会停留在 kernel 层面——每个 kernel 优化都要回到模型结构和系统成本来解释意义。

第三，\textbf{把 trade-off 显式化。}AI Infra 中几乎没有"免费的优化"。降低显存往往增加通信或计算；提高吞吐往往牺牲延迟或稳定性；简化调度往往损失灵活性。本书在每个关键选择处，都会明确写出代价和适用条件，而不是给出一个"最佳配置"就结束。

\section*{本书的结构}

全书分为五篇，外加附录：

\begin{itemize}
\item \textbf{第一篇：基石}——线性代数、微积分、概率论与 CUDA C++ 基础、深度学习基本原理。这是整本书的地基：从 shape 推导显存，从矩阵乘法推导 FLOPs，从链式法则推导反向传播与激活保存，从 softmax 推导数值稳定性，从 GPU 编程模型推导 kernel 性能。
\item \textbf{第二篇：演进}——Transformer 架构深度剖析、大语言模型发展史与主流范式、GPT/LLaMA 与现代 Decoder-only 架构。这是把地基上的数学对象组装成真实模型：Attention、FFN、RoPE、RMSNorm、SwiGLU、GQA、MoE，以及它们对训练和推理系统的具体影响。
\item \textbf{第三篇：分布式训练架构}——数据并行、张量并行、流水线并行、3D 并行策略、ZeRO 与 DeepSpeed。这是回答"大模型放不下、训不动"的核心篇章：从单卡瓶颈出发，逐步推导每种并行策略的通信模式、显存收益和适用边界，最终给出 3D 并行配置的推导方法。
\item \textbf{第四篇：推理架构优化}——KV Cache、PagedAttention 与 vLLM、FlashAttention、量化、Continuous Batching。这是回答"如何让模型高吞吐、低延迟地服务"的核心篇章：从自回归推理的 prefill/decode 两阶段出发，逐步分析 KV Cache 管理、内存碎片、IO 瓶颈、精度压缩和动态调度。
\item \textbf{第五篇：集群网络与硬件调度}——GPU 架构、NVLink/NVSwitch、RDMA/RoCE、K8s AI 调度。这是回到物理世界：理解 GPU 芯片内部结构、节点内互联、节点间网络和集群级调度，让前面的分布式策略和推理优化落到真实硬件上。
\end{itemize}

附录提供符号表、CUDA 环境配置指南、集群部署 checklist 和参考文献。

\section*{读者对象}

本书面向三类读者：

\begin{itemize}
\item \textbf{AI Infra 工程师}：日常工作中需要配置分布式训练、优化推理服务、调试 NCCL 通信、分析 GPU profiling 数据的工程师。本书帮助他们把零散经验整合为系统性理解。
\item \textbf{算法工程师转向 Infra}：已经熟悉模型结构和训练流程，但希望深入理解"为什么 OOM"、"为什么训练发散"、"为什么推理慢"的根本原因，从而更好地与 Infra 团队协作或独立解决系统问题的工程师。
\item \textbf{系统工程师转向 AI}：已经熟悉分布式系统、网络、调度和存储，但需要理解模型计算特征、显存模式和通信需求的工程师，从而把传统系统经验有效迁移到 AI 工作负载上。
\end{itemize}

无论你属于哪一类，本书都假设你具备基本的 Python 编程能力，了解深度学习的基本概念（神经网络、训练、梯度），并有过至少一次使用 GPU 运行模型的经验。不需要你预先精通线性代数、CUDA 或分布式系统——这些正是本书要从 AI Infra 角度重新教给你的。

\section*{如何使用这本书}

本书的章节之间有较强的依赖关系。第一篇是后续所有篇章的基础；第二篇是理解训练和推理篇章中模型结构的前提；第三篇和第四篇相对独立，你可以根据自己的工作重点选择先读训练篇或推理篇；第五篇为第三篇和第四篇提供硬件支撑，建议在读完至少一篇并行策略后再读。

如果你是完全的新手，建议从头到尾顺序阅读。如果你已经有分布式训练经验，可以跳过第一篇的部分章节，直接进入第三篇，但在遇到不理解的数学推导或 kernel 概念时回到第一篇补充。如果你主要做推理服务，可以先读第二篇的架构章节和第四篇，再回到第三篇了解训练侧的并行策略——因为很多推理优化（如 tensor parallel、KV Cache 管理）的训练侧对应知识同样重要。

具体的阅读路线建议，请参阅"本书阅读路线"。

\section*{致谢}

这本书的写作离不开许多人的帮助。感谢所有在开源社区中分享训练配置、推理优化、kernel 实现和调试经验的前辈和同行——你们的实践是本书最重要的素材来源。感谢 NVIDIA、PyTorch、DeepSpeed、vLLM、Megatron-LM 等团队的开源工作和文档——没有这些基础设施，AI Infra 的知识体系无法如此快速积累。感谢每一位在 GPU 集群上熬过深夜调试 loss spike、NCCL timeout 和 OOM 的工程师——你们踩过的坑，本书试图系统化地记录下来，让后来者少走弯路。

最后，感谢我的家人和朋友们在写作期间的耐心与支持。

\begin{flushright}
作者\\
\today
\end{flushright}